% section: シグマの計算

\section{シグマの計算}\label{節：シグマの計算}
  \noindent\textbf{【本編で使う場面】}\par
  和の記号と添字の役割を確認する.\sectionref*{節：階差型}で変化を足し戻し,\sectionref*{節：総和型}で総和の式をずらして引くときに使う.

  多くの項を足すとき、すべてを並べる代わりに、一つの記号でまとめられる。
  実数の\bookterm{sequence}{数列}$\{a_k\}$について、$a_1$から$a_n$までの和を
  \[
    \sum_{k=1}^{n} a_k=a_1+a_2+\cdots+ a_n
  \]
  と書く。複素数の\bookterm{sequence}{数列}でも、同じ記号を使う。
  $\sum$は、ギリシャ文字の\bookterm{sigma}{シグマ}である。
  下の$k=1$から上の$n$まで、$k$を一つずつ進めて足す。

  同じ数を繰り返し足せば、項数を掛けた値になる。
  定数倍は和の外へ出せる。二つの和は、項ごとに分けてもよい。
  \begin{align*}
    &\sum_{k=1}^{n} c=nc\\
    &\sum_{k=1}^{n} ca_k=c\sum_{k=1}^{n} a_k\\
    &\sum_{k=1}^{n} (a_k+b_k)=\sum_{k=1}^{n} a_k+\sum_{k=1}^{n} b_k
  \end{align*}
  自然数と、その二乗・三乗の和は、次の式でまとめられる。
  \begin{align*}
    &\sum_{k=1}^{n} k=\frac{1}{2}n(n+1)\\
    &\sum_{k=1}^{n} k^2=\frac{1}{6}n(n+1)(2n+1)\\
    &\sum_{k=1}^{n} k^3=\left\{\frac{1}{2}n(n+1)\right\}^2\\
  \end{align*}

  添字$k$は,和の中で順に動く番号である.同じ範囲なら$\sum_{k=1}^{n}a_k=\sum_{j=1}^{n}a_j$と書き換えても値は変わらない.
  一方,上端の$n$は足す項数を決める.例えば$S_n=\sum_{k=1}^{n}a_k$なら
  \[
    S_{n+1}=S_n+a_{n+1}
  \]
  である.
  \bookterm{recurrence}{漸化式}の添字をずらす場合は,式全体の$n$を置き換える.
  例えば$a_{n+1}=2a_n+n$から$a_{n+2}=2a_{n+1}+n+1$を得る.
